\documentclass[10pt,a4paper]{amsart}
\usepackage[margin=19mm]{geometry}
\usepackage{amsmath,amssymb,mathtools}
\usepackage[scr=rsfs,cal=euler]{mathalfa}
\usepackage{xfrac}
\usepackage[unicode,hidelinks,bookmarks=false]{hyperref}
\newcommand{\ie}{i.e.\ }
\newcommand{\cf}{cf.\ }
\newcommand{\resp}{respectively\ }
\newcommand{\Subsection}{\subsection}
\providecommand{\nobreakdash}{\nobreak\hbox{-}}
\setlength{\emergencystretch}{2em}
\multlinegap0pt
\usepackage{amssymb}
\usepackage{xcolor}
\usepackage{algorithm}\usepackage{algpseudocode}
\newtheorem{theorem}{Theorem}
\newtheorem{corollary}{Corollary}
\newcommand{\NN}{\mathbb N}
\newcommand{\ZZ}{\mathbb Z}
\title{Gap estimates for the spectrum of $m$-bonacci numbers}
\author{Anna Chiara Lai, Paola Loreti}
\date{Source: arXiv:2503.11322v2 (CC BY 4.0)}
\begin{document}
\maketitle

\noindent\emph{Source and licence.} Gap estimates for the spectrum of $m$-bonacci numbers. Version arXiv:2503.11322v2 (CC BY 4.0), distributed by arXiv under the Creative Commons Attribution 4.0 International licence, http://creativecommons.org/licenses/by/4.0/, as recorded in the arXiv licence metadata for this exact version and shown on the abstract page https://arxiv.org/abs/2503.11322v2. Full licence notice: \texttt{licenses/arxiv-2503.11322-CC-BY-4.0.txt}.\par\smallskip

\noindent\textit{Editorial scope.} Counted unit: the article's corollary cor1 (source0.tex lines 1062--1071). The article's complete original proof of it is delivered with it (source0.tex lines 1072--1116, one \texttt{proof} environment of 45 lines). No mathematics is written, repaired or re-derived: the statement and the proof are the article's own lines, in the article's own notation, and no retained line is substituted or re-flowed.  Retained uncounted as the source-local context the proof needs: lines 92--125.  Nothing was omitted from any retained span, so the package carries no omission record.  No retained line contains a citation, so the wrapper carries no bibliography and no bibliography entry.  No retained line contains a figure environment, an image-inclusion command or drawing code, so no image is delivered and the package declares no asset dependency.  Editorial formatting only, in addition: an AMS \texttt{amsart} preamble that restates the article's own declarations and macros used by the retained lines, namely the environments \texttt{theorem}, \texttt{corollary} on separate counters, together with the standard packages those lines need.  The retained environments print in this document as Theorem 1, Corollary 1 in order of appearance, whereas the article prints the same items with the numbering of its own full document.  The complete native source of the article as distributed in the arXiv e-print for this exact version is bundled unchanged as \texttt{sources/174754/source0.tex}.
\begin{theorem}\label{thm1i}
Let $m\geq 2$ and $N\in \mathbb{N}$. Then there exists a constant
$\gamma_{m,N}>0$ such that
\begin{equation}\label{gamma}
\lambda_{n+N}(q_m)-\lambda_n(q_m)\geq N\,\gamma_{m,N}
\quad \text{for all } n\geq 1.
\end{equation}

The following choice of $\gamma_{m,N}$ satisfies the above inequality.
Let $(x_h)\in\{0,1\}^\infty$ be the canonical $m$-bonacci expansion of $N$, that is,
a binary sequence with no occurrences of the block $1^m$ that expands $N$ in the form
\[
N=\sum_{h\geq 0} x_h F^{(m)}_h,
\quad x_h\in\{0,1\},
\]
where $(F^{(m)}_h)$ denotes the $m$-bonacci sequence.

Then
\begin{equation}\label{eqgamma}
\gamma_{m,N}
=\frac{1}{N}
\sum_{j=1}^m
\left(
\sum_{h\geq 0} x_h F^{(m)}_{h-j}
- b_m
\right)
d_m(j),
\end{equation}
where:
\begin{itemize}
\item[-] $d_m(j):=q_m^{j-1}-\sum_{k=0}^{j-2} q_m^k$ with $j=1,\dots,m$ are the finitely many possible gap values in $\Lambda(q_m)$;
\item[-] $b_m$ is a constant lower or equal than $m-1$ measuring the balance of the $m$-bonacci word.
\end{itemize}
\end{theorem}

\begin{corollary}[Fibonacci case]\label{cor1}
For all $N\in \NN$ we have the gap condition
\begin{equation}\label{gengapfib}
\lambda_{2,N+k}-\lambda_{2,k}\geq N\gamma_{2,N} \quad \text{for all } k\in \ZZ,
\end{equation}
where
$$\gamma_{2,N}:=\frac{1}{\varphi}+\frac{1}{\varphi^3} -\frac{1}{N}\left(\frac{1}{\varphi^4}+\varphi\right).$$
The estimate is sharpened in the case $N=2$ by setting  $$\gamma_{2,2}=\frac{\varphi}{2}.$$

\end{corollary}

\begin{proof}
Set for brevity $F^{(2)}_n=F_n$ (so that $F_0=1, F_1=2,F_2=3\cdots$, in particular $F_n$ is the $n+2$th Fibonacci number) and let $(x_h)$ be the canonical Fibonacci expansion of $N$. Note that $d_2(1)=1$ and $d_2(2)=1/\varphi$. Since the balance constant $b_2$ is equal to $1$, by Theorem \ref{thm1i} we have that 
\eqref{gengapfib} holds with 
\begin{align*}\gamma_{2,N}
&=
\frac{1}{N}
\sum_{j=1}^2
\left(
\sum_{h=0}^\infty x_h F_{h-j} - 1
\right)d_2(j)\\
&=
\frac{1}{N}
\left[
\left(\sum_{h=0}^\infty x_h F_{h-1} -1\right)
+
\left(\sum_{h=0}^\infty x_h F_{h-2} -1\right)\frac{1}{\varphi}
\right] \\
&=
\frac{1}{N}
\left(
\sum_{h=0}^\infty x_h\left(F_{h-1}+F_{h-2}\frac{1}{\varphi}\right)
-\varphi
\right).
\end{align*}
Since 
$$F_{h-1}=\frac{1}{\varphi} F_h +  \left(-\frac{1}{\varphi}\right)^{h+2};\quad F_{h-2}=\frac{1}{\varphi^2} F_h -  \left(-\frac{1}{\varphi}\right)^{h+2} $$
% for all $h$ and $\sum_{h=0}^\infty x_h F_h=N$, 
we obtain
\begin{align*}\gamma_{2,N}&=\frac{1}{ N}\left(\left(\frac{1}{\varphi}+\frac{1}{\varphi^3}\right) \sum_{h=0}^\infty x_h F_h+ \frac{1}{\varphi^4}\sum_{h=0}^\infty  \frac{x_h}{(-\varphi)^h}-\varphi\right)\\
&=\frac{1}{\varphi}+\frac{1}{\varphi^3} +\frac{1}{ N}\left(\frac{1}{\varphi^4}\sum_{h=0}^\infty \frac{x_h}{(-\varphi)^h}-\varphi\right)\\
&\geq\frac{1}{\varphi}+\frac{1}{\varphi^3}  -\frac{1}{ N}\left(\frac{1}{\varphi^3}\sum_{h=1}^\infty \frac{1}{\varphi^{2h}}+\varphi\right)\\
&= \frac{1}{\varphi}+\frac{1}{\varphi^3} -\frac{1}{N}\left(\frac{1}{\varphi^4}+\varphi\right).
\end{align*}


In the particular case $N=2$, we note that the subword $22$ does not occur
in the Fibonacci word.
Therefore,
\[
\lambda_{2,k+2}-\lambda_{2,k}\ge d_2(1)+d_2(2)=\varphi,
\]
and one can choose $\gamma_{2,2}=\varphi/2$.


\end{proof}

\end{document}
